\documentclass[UTF8,11pt]{ctexart}
\usepackage[a4paper,margin=24mm]{geometry}
\usepackage{amsmath,amssymb,amsthm,mathtools,mathrsfs}
\usepackage[all]{xy}
\usepackage{xurl,hyperref,enumitem}
\usepackage{tikz}
\usetikzlibrary{arrows.meta,cd,decorations.markings,calc,patterns}
\hypersetup{hidelinks,breaklinks=true}
\setlength{\emergencystretch}{3em}
\allowdisplaybreaks
\newtheorem*{theorem}{Theorem}
\newtheorem*{thm}{Theorem}
\newtheorem*{lemma}{Lemma}
\newtheorem*{proposition}{Proposition}
\newtheorem*{prop}{Proposition}
\newtheorem*{corollary}{Corollary}
\newtheorem*{cor}{Corollary}
\newtheorem*{claim}{Claim}
\theoremstyle{definition}
\newtheorem*{definition}{Definition}
\newtheorem*{defn}{Definition}
\newtheorem*{remark}{Remark}
\newtheorem*{example}{Example}
\newcommand{\R}{\mathbb R}
\newcommand{\C}{\mathbb C}
\newcommand{\Z}{\mathbb Z}
\newcommand{\Q}{\mathbb Q}
\newcommand{\N}{\mathbb N}
\newcommand{\F}{\mathbb F}
\newcommand{\D}{\mathbb D}
\newcommand{\bB}{\mathbb B}
\newcommand{\bC}{\mathbb C}
\newcommand{\bR}{\mathbb R}
\newcommand{\bZ}{\mathbb Z}
\newcommand{\bN}{\mathbb N}
\newcommand{\bQ}{\mathbb Q}
\newcommand{\bD}{\mathbb D}
\newcommand{\bF}{\mathbb F}
\newcommand{\CP}{\mathbb{CP}}
\newcommand{\RP}{\mathbb{RP}}
\newcommand{\sO}{\mathscr O}
\newcommand{\sI}{\mathscr I}
\newcommand{\sF}{\mathscr F}
\newcommand{\sG}{\mathscr G}
\newcommand{\sH}{\mathscr H}
\newcommand{\sR}{\mathscr R}
\newcommand{\sM}{\mathscr M}
\newcommand{\sV}{\mathscr V}
\newcommand{\sU}{\mathscr U}
\DeclareMathOperator{\Hom}{Hom}
\DeclareMathOperator{\Ext}{Ext}
\DeclareMathOperator{\Tor}{Tor}
\DeclareMathOperator{\Spec}{Spec}
\DeclareMathOperator{\Proj}{Proj}
\DeclareMathOperator{\supp}{supp}
\DeclareMathOperator{\im}{im}
\DeclareMathOperator{\id}{id}
\DeclareMathOperator{\Tr}{Tr}
\DeclareMathOperator{\tr}{tr}
\DeclareMathOperator{\rank}{rank}
\DeclareMathOperator{\ord}{ord}
\DeclareMathOperator{\dist}{dist}
\DeclareMathOperator{\End}{End}
\DeclareMathOperator{\Aut}{Aut}
\DeclareMathOperator{\Gal}{Gal}
\DeclareMathOperator{\vol}{vol}
\DeclareMathOperator{\Cl}{Cl}
\DeclareMathOperator{\coker}{coker}
\DeclareMathOperator{\codim}{codim}
\DeclareMathOperator{\divergence}{div}
\DeclareMathOperator{\Ric}{Ric}
\DeclareMathOperator{\grad}{grad}
\DeclareMathOperator{\Hess}{Hess}
\newcommand{\abs}[1]{\left\lvert#1\right\rvert}
\newcommand{\sabs}[1]{\lvert#1\rvert}
\newcommand{\babs}[1]{\bigl\lvert#1\bigr\rvert}
\newcommand{\Babs}[1]{\Bigl\lvert#1\Bigr\rvert}
\newcommand{\bbabs}[1]{\biggl\lvert#1\biggr\rvert}
\newcommand{\BBabs}[1]{\Biggl\lvert#1\Biggr\rvert}
\newcommand{\norm}[1]{\left\lVert#1\right\rVert}
\newcommand{\snorm}[1]{\lVert#1\rVert}
\newcommand{\bnorm}[1]{\bigl\lVert#1\bigr\rVert}
\newcommand{\Bnorm}[1]{\Bigl\lVert#1\Bigr\rVert}
\newcommand{\bbnorm}[1]{\biggl\lVert#1\biggr\rVert}
\newcommand{\BBnorm}[1]{\Biggl\lVert#1\Biggr\rVert}
\newcommand{\widebar}[1]{\overline{#1}}
\newcommand{\nicefrac}[2]{\frac{#1}{#2}}
\newcommand{\linnprod}[2]{\langle#1,#2\rangle}
\newcommand{\myindex}[1]{#1}
\newcommand{\glsadd}[1]{}
\newcommand{\avoidbreak}{}
\newcommand{\sourceref}[1]{\textnormal{[原文：\nolinkurl{#1}]}}
\newcommand{\thmref}[1]{\sourceref{#1}}
\newcommand{\lemmaref}[1]{\sourceref{#1}}
\newcommand{\propref}[1]{\sourceref{#1}}
\newcommand{\corref}[1]{\sourceref{#1}}
\newcommand{\defnref}[1]{\sourceref{#1}}
\newcommand{\exampleref}[1]{\sourceref{#1}}
\newcommand{\exerciseref}[1]{\sourceref{#1}}
\newcommand{\figureref}[1]{\sourceref{#1}}
\newcommand{\sectionref}[1]{\sourceref{#1}}
\newcommand{\appendixref}[1]{\sourceref{#1}}
\newcommand{\stacksref}[2]{\href{https://stacks.math.columbia.edu/tag/#1}{#2 [#1]}}


\newcommand{\E}{\mathbb E}
\newcommand{\Prob}{\mathbb P}
\newcommand{\Reff}{\mathcal R_{\mathrm{eff}}}
\newcommand{\eps}{\varepsilon}
\DeclareMathOperator{\diam}{diam}
\DeclareMathOperator{\Var}{Var}
\DeclareMathOperator{\Crit}{Crit}
\newcommand{\dd}{\,\mathrm d}

\begin{document}
\section*{071\quad 三元有限域无直线集的切片秩上界}
\subsection*{来源与计数目标}
Terence Tao, \emph{A symmetric formulation of the Croot--Lev--Pach--Ellenberg--Gijswijt capset bound}, 18 May 2016, What\textquotesingle s new。采用网页实际访问正文（2026-09-28），公式 (1)--(3)、Lemmas 1--2 及其后 $|A|\le3N$ 的推论。
\par\url{https://terrytao.wordpress.com/2016/05/18/a-symmetric-formulation-of-the-croot-lev-pach-ellenberg-gijswijt-capset-bound/}
\par\url{https://terrytao.wordpress.com/about/}
\par\textbf{本文件只计一个问题（整理者概述）：}对 $A\subset\mathbb F_3^n$ 不含非平凡直线，证明 $|A|\le3\sum_{a+b+c=n,\ b+2c\le2n/3}n!/(a!b!c!)$；对角超矩阵秩引理并入本题，不另计数。
\subsection*{作者命题与人类证明}

A capset in the vector space ${\bf F}_3^n$ over the finite field ${\bf F}_3$ of three elements is a subset $A$ of ${\bf F}_3^n$ that does not contain any lines $\{x,x+r,x+2r\}$, where $x,r\in{\bf F}_3^n$ and $r\ne0$.

The basic starting point is this: if $A$ is a capset, then one has the identity
\begin{equation}
\delta_{0^n}(x+y+z)=\sum_{a\in A}\delta_a(x)\delta_a(y)\delta_a(z)\tag{1}
\end{equation}
for all $(x,y,z)\in A^3$, where $\delta_a(x):=1_{a=x}$ is the Kronecker delta function, which we view as taking values in ${\bf F}_3$. Indeed, (1) reflects the fact that the equation $x+y+z=0$ has solutions precisely when $x,y,z$ are either all equal, or form a line, and the latter is ruled out precisely when $A$ is a capset.

To exploit (1), we will show that the left-hand side of (1) is ``low rank'' in some sense, while the right-hand side is ``high rank''.

Recall that a function $F:A\times A\to{\bf F}$ taking values in a field ${\bf F}$ is of rank one if it is non-zero and of the form $(x,y)\mapsto f(x)g(y)$ for some $f,g:A\to{\bf F}$, and that the rank of a general function $F:A\times A\to{\bf F}$ is the least number of rank one functions needed to express $F$ as a linear combination.

More generally, if $k\ge2$, we define the rank of a function $F:A^k\to{\bf F}$ to be the least number of ``rank one'' functions of the form
\[
(x_1,\ldots,x_k)\mapsto f(x_i)g(x_1,\ldots,x_{i-1},x_{i+1},\ldots,x_k)
\]
for some $i=1,\ldots,k$ and some functions $f:A\to{\bf F}$, $g:A^{k-1}\to{\bf F}$, that are needed to generate $F$ as a linear combination. For instance, when $k=3$, the rank one functions take the form $(x,y,z)\mapsto f(x)g(y,z)$, $(x,y,z)\mapsto f(y)g(x,z)$, $(x,y,z)\mapsto f(z)g(x,y)$, and linear combinations of $r$ such rank one functions will give a function of rank at most $r$.

It is a standard fact in linear algebra that the rank of a diagonal matrix is equal to the number of non-zero entries. This phenomenon extends to higher dimensions:
\begin{lemma}[1, Rank of diagonal hypermatrices]
Let $k\ge2$, let $A$ be a finite set, let ${\bf F}$ be a field, and for each $a\in A$, let $c_a\in{\bf F}$ be a coefficient. Then the rank of the function
\begin{equation}
(x_1,\ldots,x_k)\mapsto\sum_{a\in A}c_a\delta_a(x_1)\cdots\delta_a(x_k)\tag{2}
\end{equation}
is equal to the number of non-zero coefficients $c_a$.
\end{lemma}
\begin{proof}
We induct on $k$. As mentioned above, the case $k=2$ follows from standard linear algebra, so suppose now that $k>2$ and the claim has already been proven for $k-1$.

It is clear that the function (2) has rank at most equal to the number of non-zero $c_a$ (since the summands on the right-hand side are rank one functions), so it suffices to establish the lower bound. By deleting from $A$ those elements $a\in A$ with $c_a=0$ (which cannot increase the rank), we may assume without loss of generality that all the $c_a$ are non-zero. Now suppose for contradiction that (2) has rank at most $|A|-1$, then we obtain a representation
\begin{equation}
\begin{split}
&\sum_{a\in A}c_a\delta_a(x_1)\cdots\delta_a(x_k)\\
&\quad=\sum_{i=1}^k\sum_{\alpha\in I_i}f_{i,\alpha}(x_i)g_{i,\alpha}(x_1,\ldots,x_{i-1},x_{i+1},\ldots,x_k)
\end{split}\tag{3}
\end{equation}
for some sets $I_1,\ldots,I_k$ of cardinalities adding up to at most $|A|-1$, and some functions $f_{i,\alpha}:A\to{\bf F}$ and $g_{i,\alpha}:A^{k-1}\to{\bf R}$.

Consider the space of functions $h:A\to{\bf F}$ that are orthogonal to all the $f_{k,\alpha}$, $\alpha\in I_k$ in the sense that
\[
\sum_{x\in A}f_{k,\alpha}(x)h(x)=0
\]
for all $\alpha\in I_k$. This space is a vector space whose dimension $d$ is at least $|A|-|I_k|$. A basis of this space generates a $d\times|A|$ coordinate matrix of full rank, which implies that there is at least one non-singular $d\times d$ minor. This implies that there exists a function $h:A\to{\bf F}$ in this space which is nowhere vanishing on some subset $A'$ of $A$ of cardinality at least $|A|-|I_k|$.

If we multiply (3) by $h(x_k)$ and sum in $x_k$, we conclude that
\[
\begin{split}
&\sum_{a\in A}c_ah(a)\delta_a(x_1)\cdots\delta_a(x_{k-1})\\
&\quad=\sum_{i=1}^{k-1}\sum_{\alpha\in I_i}f_{i,\alpha}(x_i)\tilde g_{i,\alpha}(x_1,\ldots,x_{i-1},x_{i+1},\ldots,x_{k-1})
\end{split}
\]
where
\[
\begin{split}
&\tilde g_{i,\alpha}(x_1,\ldots,x_{i-1},x_{i+1},\ldots,x_{k-1})\\
&\quad:=\sum_{x_k\in A}g_{i,\alpha}(x_1,\ldots,x_{i-1},x_{i+1},\ldots,x_k)h(x_k).
\end{split}
\]
The right-hand side has rank at most $|A|-1-|I_k|$, since the summands are rank one functions. On the other hand, from induction hypothesis the left-hand side has rank at least $|A|-|I_k|$, giving the required contradiction.
\end{proof}

On the other hand, we have the following (symmetrised version of a) beautifully simple observation of Croot, Lev, and Pach:
\begin{lemma}[2]
On $({\bf F}_3^n)^3$, the rank of the function $(x,y,z)\mapsto\delta_{0^n}(x+y+z)$ is at most $3N$, where
\[
N:=\sum_{\substack{a,b,c\ge0:\ a+b+c=n\ b+2c\le2n/3}}\frac{n!}{a!b!c!}.
\]
\end{lemma}
\begin{proof}
Using the identity $\delta_0(x)=1-x^2$ for $x\in{\bf F}_3$, we have
\[
\delta_{0^n}(x+y+z)=\prod_{i=1}^n(1-(x_i+y_i+z_i)^2).
\]
The right-hand side is clearly a polynomial of degree $2n$ in $x,y,z$, which is then a linear combination of monomials
\[
x_1^{i_1}\cdots x_n^{i_n}y_1^{j_1}\cdots y_n^{j_n}z_1^{k_1}\cdots z_n^{k_n}
\]
with $i_1,\ldots,i_n,j_1,\ldots,j_n,k_1,\ldots,k_n\in\{0,1,2\}$ with
\[
i_1+\cdots+i_n+j_1+\cdots+j_n+k_1+\cdots+k_n\le2n.
\]
In particular, from the pigeonhole principle, at least one of $i_1+\cdots+i_n,j_1+\cdots+j_n,k_1+\cdots+k_n$ is at most $2n/3$.

Consider the contribution of the monomials for which $i_1+\cdots+i_n\le2n/3$. We can regroup this contribution as
\[
\sum_\alpha f_\alpha(x)g_\alpha(y,z)
\]
where $\alpha$ ranges over those $(i_1,\ldots,i_n)\in\{0,1,2\}^n$ with $i_1+\cdots+i_n\le2n/3$, $f_\alpha$ is the monomial
\[
f_\alpha(x_1,\ldots,x_n):=x_1^{i_1}\cdots x_n^{i_n}
\]
and $g_\alpha:{\bf F}_3^n\times{\bf F}_3^n\to{\bf F}_3$ is some explicitly computable function whose exact form will not be of relevance to our argument. The number of such $\alpha$ is equal to $N$, so this contribution has rank at most $N$. The remaining contributions arising from the cases $j_1+\cdots+j_n\le2n/3$ and $k_1+\cdots+k_n\le2n/3$ similarly have rank at most $N$ (grouping the monomials so that each monomial is only counted once), so the claim follows.
\end{proof}

Upon restricting from $({\bf F}_3^n)^3$ to $A^3$, the rank of $(x,y,z)\mapsto\delta_{0^n}(x+y+z)$ is still at most $3N$. The two lemmas then combine to give the Ellenberg--Gijswijt bound
\[
|A|\le3N.
\]

\subsection*{整理说明与先修范围（非作者正文）}
保留作者定义的 ``rank''（当前通常称切片秩）及证明次序。计数目标为正文已推出的精确有限和上界，不额外要求后面的数值渐近常数；没有收录后续 Remarks 3--6。Lemmas 1--2 的全部论证和主目标返回步骤均在正文。

允许先修为有限维线性代数、矩阵秩和有限域基本算术；本题仍需独立构造三变量指示张量，证明对角切片秩下界，并把多项式总次数转为切片数量约束。并非直接引用 capset 界。与007的方向覆盖问题、008的直线交会问题不同，所用低秩证据也不同。

原网页公式(3)后 $g_{i,\alpha}$ 的值域印成 ${\bf R}$，其上文定义值域为 ${\bf F}$；此处保留网页字样并记录，不静默改写作者正文。公式分行和链接编号作独立排版。
\subsection*{可修改的简短标注（非作者正文）}
$c$：有限域无三项等差结构点集的极值大小。

$a$：指示张量；切片秩；对角超矩阵；正交补与非零子式；有限域多项式；总次数分配；鸽巢原理。

\texttt{cross\_theory\_status = yes}。点集的禁线条件把加法约束张量变为对角张量；多项式展开控制其切片数，从而将代数秩比较返回点集大小。位置：公式(1)、Lemmas 1--2及随后主界。

全部标注为非穷尽、可修改初稿。
\subsection*{来源许可与修改记录}
作者 About 页 ``Copyright etc. 明确允许复制、引用或翻译合理部分（例如单篇文章），较长转载需注明原URL。本题依该明示许可摘取一篇文章的相关证明；保留作者和来源，不将许可改称CC。
2026-09-28：ChatGPT 整理为独立 TeX，加入编号、来源定位、整理说明与初稿标注；原作者数学论证与整理者文字分开标示。
\end{document}
